\\chapter*{Introduction générale}
\\addcontentsline{toc}{chapter}{Introduction générale} % to include the introduction to the table of content

La transformation digitale constitue aujourd'hui un levier stratégique majeur pour les entreprises industrielles. Dans un environnement économique de plus en plus compétitif, les organisations doivent s'appuyer sur des systèmes d'information performants et fiables pour maintenir leur avantage concurrentiel. Le parc informatique, composé de l'ensemble des ressources matérielles et logicielles, représente un pilier fondamental de cette infrastructure numérique.

Depuis l'avènement de l'informatique, de nombreuses activités de la vie quotidienne ont été simplifiées. Les logiciels et les réseaux informatiques ont permis aux individus de traiter et de gérer facilement des informations. De nos jours, ces systèmes informatiques sont omniprésents dans la plupart des grandes entreprises, quel que soit leur secteur d'activité. Dans le secteur de la fabrication électronique (EMS -- Electronic Manufacturing Services), où la précision et la traçabilité sont des exigences fondamentales, la gestion efficace du parc informatique revêt une importance particulière.

Avec l'essor d'Internet et de l'industrie 4.0, de plus en plus d'entreprises ouvrent leur parc informatique à leurs partenaires et fournisseurs. Le parc informatique désigne l'ensemble des ressources matérielles et logicielles utilisées par une société pour le traitement automatisé de l'information. Pour assurer la survie, la gestion et la pérennité de ces ressources, il est essentiel d'avoir une gestion efficace du parc informatique de l'entreprise. Cette gestion, souvent désignée sous le terme ITSM (IT Service Management), s'inscrit dans le cadre des bonnes pratiques ITIL (Information Technology Infrastructure Library).

La gestion du parc informatique comprend plusieurs aspects essentiels. Tout d'abord, il s'agit de recenser et de localiser tous les équipements informatiques de l'entreprise, tels que les ordinateurs de bureau, les ordinateurs portables, les serveurs, les switchs réseau, les imprimantes et autres périphériques. Il est important de connaître leur emplacement, leur état, leur configuration et leur historique afin de pouvoir les gérer efficacement et d'optimiser leur cycle de vie.

Ensuite, la gestion du parc informatique implique la réalisation de tâches de maintenance et d'assistance aux utilisateurs. Cela inclut la mise à jour des logiciels, le suivi des licences, la résolution des problèmes techniques, le support aux utilisateurs en cas de difficultés et la planification des renouvellements d'équipements. La traçabilité des affectations et des mouvements de matériel constitue également un enjeu majeur pour garantir la sécurité et la conformité des actifs informatiques.

Cependant, ces opérations peuvent souvent dépasser les compétences d'une seule personne et devenir difficiles à gérer manuellement, notamment lorsque le parc informatique atteint une taille conséquente. C'est pourquoi il est nécessaire de mettre en place des outils et des systèmes au sein de l'entreprise pour assurer un suivi régulier du parc informatique, automatiser les tâches récurrentes et anticiper les éventuelles défaillances.

Effectivement, dans le cadre de notre projet de fin d'études, nous avons été mandatés par la société ASTEELFLASH, un acteur majeur de l'industrie EMS (Electronic Manufacturing Services) en Tunisie, pour développer une application web de gestion du parc informatique et CMDB (Configuration Management Database). ASTEELFLASH, avec ses 1000 employés et ses 15 000 m² de surface de production à La Soukra, gère un parc informatique conséquent composé de centaines d'équipements de différents types, ce qui justifie pleinement le besoin d'un outil de gestion dédié.

L'objectif de cette application est de fournir à ASTEELFLASH une solution complète et efficace pour gérer l'ensemble de son parc informatique. Cela comprend la possibilité de recenser et de localiser tous les équipements informatiques, de suivre leur état, leur configuration et leurs caractéristiques. De plus, l'application permettra de planifier et de réaliser les opérations de maintenance, de gérer les affectations d'équipements aux employés, de suivre les missions temporaires de matériel et de recevoir des alertes automatiques en cas de stock critique.

Le présent rapport est structuré en cinq chapitres, décrivant les différentes étapes de notre travail :

\\begin{itemize}
  
\\item \\textbf{Chapitre 1 -- Cadre général du projet :} Ce chapitre présente le contexte général du projet en introduisant l'organisme d'accueil ASTEELFLASH, ses activités et son organigramme. Il aborde ensuite l'étude de l'existant, la problématique identifiée et la solution proposée. Une étude comparative des solutions existantes sur le marché (GLPI, OCS Inventory, Snipe-IT, Lansweeper) justifie le choix du développement sur mesure. Enfin, le chapitre présente les concepts ITIL/CMDB et la méthodologie agile Scrum adoptée.

\\item \\textbf{Chapitre 2 -- Analyse et spécification des besoins :} Ce chapitre traite de la spécification des besoins fonctionnels et non fonctionnels du système. Il comprend le diagramme de cas d'utilisation global, le diagramme de classes, la planification Scrum et le backlog produit détaillé.

\\item \\textbf{Chapitre 3 -- Release 1 : Architecture et Sprint 1 :} Ce chapitre se concentre sur les choix technologiques (ASP.NET Core, Entity Framework Core, SQL Server), l'architecture 3-tiers et MVC de l'application, le modèle de données, les mécanismes de sécurité LDAP, le sprint 0 (mise en place du socle) et le sprint 1 dédié à l'authentification et à la gestion des profils.

\\item \\textbf{Chapitre 4 -- Release 2 : Sprints 2 et 3 :} Ce chapitre couvre la gestion des équipements, le suivi statistique via le tableau de bord, la gestion des affectations, les missions temporaires d'équipements avec rappels par e-mail, et les alertes de stock minimum. Chaque sprint est présenté avec son backlog, ses diagrammes UML et ses interfaces réalisées.

\\item \\textbf{Chapitre 5 -- Release 3 : Sprint 4 et Tests :} Ce chapitre concerne la gestion des commandes, des fournisseurs et des employés. Il inclut également une section dédiée aux tests et à la validation de l'application, avec les résultats des tests fonctionnels et de performance.
 
\\end{itemize}
Ce rapport se termine par une conclusion générale présentant un bilan du travail réalisé, les compétences acquises et les perspectives d'évolution pour d'éventuelles améliorations futures.



